% Lösungen Einheit 2 von 2 – Vorgegebene Regeln anwenden (zusammenbau v0.1)
\einheitenkopf{Einheit 2 von 2 · Vorgegebene Regeln anwenden}

\erg{7}{a) der Exponent \quad b) $g(x) = x^2 + 3$, $g'(x) = 2x$ \quad c) $g(x) = 1 - x^2$, $g'(x) = -2x$, also $f(x) = -g'(x) \cdot e^{g(x)}$ \quad d) $\left[e^{x^2}\right]_0^1 = e - 1 \approx 1{,}718$ \quad e) $g(x) = 2 - x^2$, $g'(x) = -2x$, also $f = -g' \cdot e^g$: $-\left[e^{2 - x^2}\right]_0^1 = e^2 - e \approx 4{,}671$}
\erg{8}{a) Mia übersieht das Vorzeichen: $g'(x) = -2x$, also $-\left[e^{-x^2}\right]_0^1 = 1 - e^{-1} \approx 0{,}632$ \quad b) Sie ist die Kettenregel rückwärts: Die Ableitung von $e^{g(x)}$ ist genau $g'(x) \cdot e^{g(x)}$; fehlt der Faktor g', ist $e^g$ keine Stammfunktion.}
